\section{Carrying the phase through a pipeline}
\label{sec:pipeline}

This section fixes where each component of $\can(P)=(\hat P,\epsilon,k)$
is computed, stored and consumed in the three implementations the
benchmarks are drawn from, so that the numbers in
Section~\ref{sec:benchmarks} can be read against a definite data
flow. The three are (i) the Rust index-calculus pipeline with a folded
pair table and the signed-Frobenius Pollard rho reference it is
measured against, both in the sibling repository
(\sibling, degrees $n\le 63$ in one machine word);
(ii) the $n=83$ packed root index of Section~\ref{sec:record}
(\thislib, two-limb words); and (iii) the Sage pipeline for the $n=19$
isogeny-class study, which is the one that produces verified
logarithms with a $\Gamma$-invariant base and a Gr\"obner
decomposition.

\subsection{One convention per pipeline}
\label{sec:conventions}

Four canonical-representative conventions coexist in the code, and
they agree on \emph{which} points form an orbit but not on \emph{which}
member represents it:
\begin{enumerate}[leftmargin=1.6em,itemsep=0.2em]
\item \emph{First met.} The factor-base builder visits
  $\pm\tau^k(P)$ for $k=0,\dots,n-1$ from each unassigned point and
  records $(\text{orbit},k,\text{negated})$ on first sight. Cheap,
  order-dependent, and used only where the base is enumerated
  once.\art{crypto:src/cryptanalysis/koblitz_index_calculus.rs, orbit_maps_packed}
\item \emph{Lexicographic minimum over the $2n$ members} of
  $\{\pm\tau^k([h]P)\}$ by the key $(x+1,y)$, after cofactor
  projection. Deterministic and basis-independent; costs $4\,n$
  coordinate squarings per column to build. Used for the relation
  columns.\art{crypto:src/cryptanalysis/koblitz_index_calculus.rs, build_projected_signed_orbit_map}
\item \emph{Least rotation in a normal basis} of the $x$ word, with the
  sign chosen afterwards by comparing $y$ with $y+x$ (in normal
  coordinates in the rho walk, in polynomial coordinates in the
  boundary harness). This is Definition~\ref{def:sf} with $\prec$ the
  numeric order of the normal-basis word, and it is the convention
  every per-probe canonicalization uses.\art{crypto:src/cryptanalysis/koblitz_fast.rs, FrobeniusCanon; crypto:src/cryptanalysis/koblitz_strong_rho.rs, canonicalize}
\item \emph{Least numeric rotation in the permuted type-II ONB} of
  the ECC2K-130 client, where $\tau$ is a coordinate permutation
  rather than a rotation, and a different member of the orbit is
  chosen than by (3).\art{crypto:ecc2k130/CAIRN-WITNESS.md}
\end{enumerate}
The normal element in (3) is found by a bounded randomized search, so
two independent builds name orbits by different, equally valid,
incompatible keys. A folded table is therefore not portable across a
basis change, and the GPU kernels import the host's basis-change
tables rather than deriving their own.\art{crypto:src/cryptanalysis/koblitz_fast.rs, FrobeniusCanon::new} We record this as a
requirement for any reproducible artefact: \emph{the normal element,
or the basis-change tables, are part of the factor-base recipe and
must be archived with it.}

\subsection{The index-calculus data flow (\sibling, $n\le63$)}

\begin{enumerate}[leftmargin=1.6em,itemsep=0.2em]
\item \emph{Factor base.} Points with $x$ in a chosen set, closed
  under $\Gamma$; columns are signed orbits after cofactor
  projection (convention~2); a point with $[h]P=O$ gets no column.
\item \emph{Pair table.} One of three tiers
  (Section~\ref{sec:pairtable}); the folded tier keys by
  convention~3 on $x(P_i+P_j)$ and tags the row orbit.
\item \emph{Collection.} For each random $R$, the $m=3$ scan walks
  the base, forms $R-P_k$, canonicalizes its $x$ (convention~3,
  blocked $1024$ at a time) and probes the table. A hit yields a
  candidate pair; the orbit tag turns recovery into a walk over $2n$
  points; the candidate relation is checked as a point identity
  before admission.
\item \emph{Phase to coefficient.} Each admitted summand is mapped to
  its column by convention~2, giving $(c,k,\text{negated})$, and the
  coefficient $\pm\lambda^{k}$ is accumulated in column $c$
  (Proposition~\ref{prop:transport}). Two summands in one column add.
\item \emph{Linear algebra and recovery.} Sparse elimination modulo
  $r$; the target's own relation gives $s$; every recovered logarithm
  is verified as $[\log]P_0=P$.
\end{enumerate}
The phase appears twice (steps 3 and 4) and is never stored: the
table stores the tag, the matrix stores $\lambda^k$.

\subsection{The matched reference: signed-Frobenius Pollard rho}
\label{sec:rho-ref}

The reference against which every index-calculus number in this paper
is read is Pollard rho on classes of $\Gamma$, the Wiener--Zuccherato
and Gallant--Lambert--Vanstone walk with a $\sqrt{2n}$ saving
\citep{wienerzuccherato1998,glv2001}. Its iteration function has to
be class-invariant, and the implementation makes it so by the same
canonicalization as the index-calculus side: least rotation of the
$x$ word in a normal basis, sign by comparing $y$ with $y+x$ in
normal coordinates, orbit multiplier $\pm\lambda^k$ from a table,
with an exhaustive fallback over the $2n$ members when the word is
constant.\art{crypto:src/cryptanalysis/koblitz_strong_rho.rs, canonicalize} So the two arms of every comparison pay the same canonicalization
for the same automorphisms, which is the condition under which the
comparison measures the index calculus and not the representation.
Section~\ref{sec:bench-rho} shows how much of the reference's own cost
the canonicalization is, which is the reason this condition matters.

\subsection{The $n=83$ root index (\thislib)}

The decomposition is a meet-in-the-middle on roots of $S_3$: for
sampled pairs of base points the two roots are canonicalized
(convention~3, Algorithm~\ref{alg:zerogap}) and stored with their
shift (Section~\ref{sec:record}); the target's roots are canonicalized
and probed; a hit gives the relative shift by subtraction modulo
$n$, and the candidate relation is checked natively. The published
runs stop at the frozen state cap with no relation, so this pipeline
contributes stage measurements only (Section~\ref{sec:bench-canon}).

\subsection{The $n=19$ Sage pipeline (\thislib)}
\label{sec:sage-pipeline}

The isogeny-class study runs complete index calculus on
$E_0/\FF_{2^{19}}$ with a prime subgroup of order $130\,873$ and
compares a plain subspace base, $x\in V=\langle1,z,\dots,z^{9}\rangle$,
with the $\Gamma$-closure of a smaller subspace,
$F=\{\tau^jP: x(P)\in V'\}$ with $\dim V'\in\{6,7\}$. Orbit
representatives are $\min_e\,x^{2^e}$ in the integer encoding of the
polynomial basis (a convention~2 choice, affordable at $n=19$), and
each base point is stored as $(P,c,e,\pm1)$. The decomposition
$R=P_1\pm\tau^jP_2$ is a $2\dim V'$-variable quadratic Boolean system
for each shift $j$ (squaring is $\FF_2$-linear, so $\tau^j$ on $V'$ is
a linear change of the second block of variables); the shifts are
tried in a random order and solved by PolyBoRi Gr\"obner bases. A hit
returns $[(c_1,\epsilon_1,e_1),(c_2,\epsilon_2,(j+e_2)\bmod n)]$ and
the row gets $\epsilon\lambda^{e}$ in column $c$, with $\lambda$ the
root of $\lambda^2+\lambda+2$ for which
$[\lambda]G=(x_G^2,y_G^2)$.\art{experiments/volcano-ic/tau.sage, tauRelation and tauDlp} The matrix is solved by a sparse left kernel modulo $r$, and each
of the ten planted logarithms is checked against its scalar. This is
the pipeline behind the only end-to-end, verified, phase-aware
speedup in the paper (Section~\ref{sec:bench-e2e}).
